\section{Konversi Tipe dan Pelaporan Kesalahan}
\label{sec:ch15-konversi}

\subsection{Konversi Implisit}

Konversi implisit (\eng{coercion}) terjadi ketika tipe sumber dapat
dipromosikan ke tipe target tanpa kehilangan informasi:
\begin{itemize}[leftmargin=*]
  \item \code{int} ke \code{float};
  \item \code{char} ke \code{int};
  \item \code{char} ke \code{float}.
\end{itemize}

Konversi sebaliknya memerlukan cast eksplisit. Pada Subset-C v1.0,
cast belum didukung, sehingga \code{float} ke \code{int} merupakan
kesalahan.

\subsection{Konversi Eksplisit}

Cast eksplisit ditulis \code{(int) ekspresi}. Cast memberi tahu
compiler bahwa kehilangan presisi disengaja. Jika cast belum
didukung, compiler harus melaporkan kesalahan dengan jelas, bukan
diam-diam melakukan konversi.

\begin{table}[H]
\centering
\caption{Aturan konversi Subset-C v1.0}
\label{tab:ch15-konversi}
\begin{tabularx}{\textwidth}{@{}l l l X@{}}
\toprule
\textbf{Dari} & \textbf{Ke} & \textbf{Implisit?} & \textbf{Catatan} \\
\midrule
\code{int} & \code{float} & ya & promosi \\
\code{char} & \code{int} & ya & promosi \\
\code{char} & \code{float} & ya & promosi \\
\code{float} & \code{int} & tidak & memerlukan cast \\
\code{int} & \code{char} & tidak & berpotensi kehilangan data \\
\code{void} & tipe lain & tidak & tidak valid \\
\bottomrule
\end{tabularx}
\end{table}

\subsection{Pesan Kesalahan}

Pesan kesalahan tipe harus menyebut:
\begin{itemize}[leftmargin=*]
  \item lokasi ekspresi atau deklarasi;
  \item tipe yang diharapkan;
  \item tipe yang ditemukan;
  \item saran perbaikan bila memungkinkan.
\end{itemize}

\begin{table}[H]
\centering
\caption{Contoh pesan kesalahan tipe}
\label{tab:ch15-error}
\begin{tabularx}{\textwidth}{@{}l X@{}}
\toprule
\textbf{Masukan} & \textbf{Pesan} \\
\midrule
\code{int x = 3.5;} &
\code{baris 1, kolom 9: tidak dapat menetapkan float ke int} \\
\code{int x; x + "halo";} &
\code{baris 1, kolom 12: operator aritmetika memerlukan operand numerik} \\
\code{return q;} &
\code{baris 1, kolom 8: identifier 'q' tidak dideklarasikan} \\
\code{int x; int x;} &
\code{baris 1, kolom 8: identifier 'x' sudah dideklarasikan} \\
\bottomrule
\end{tabularx}
\end{table}

\subsection{Propagasi Error}

Nilai \code{TYPE\_ERROR} meredam kesalahan turunan. Jika satu
operand error, hasil operasi juga error tanpa pesan baru. Strategi
ini menjaga laporan tetap fokus pada akar masalah.

\begin{warningbox}
Jangan menampilkan semua kesalahan sekaligus tanpa prioritas.
Kelompokkan berdasarkan akar masalah dan urutkan berdasarkan lokasi
agar pengguna dapat memperbaiki secara sistematis.
\end{warningbox}
